\chapter{Máquinas de estados}
\section{\obj}
\capacidad
\begin{itemize}
\item Crear máquinas de estados para automatizar procesos de medición.
\end{itemize}

\section{\mat}
\textbf{A suministrar por la Escuela:}
\begin{itemize}
\item 1 dispositivo de adquisición de datos NI myDAQ
\item 1 termistor TDC310 de \SI{10}{\kilo\ohm} 
\item 1 resistor de \SI{10}{\kilo\ohm}
\item 1 capacitor de \SI{0.1}{\micro\farad}

\end{itemize}
\textbf{A suministrar por el estudiante:}
\begin{itemize}
\item 1 breadboard
\item Cables de interconexión (jumpers) macho-macho
\end{itemize}

\section{\pro}
\begin{enumerate}
\item Realice las conexiones del circuito tal y como se indica en la Figura \ref{fig:L3F1}.
\item Cree un instrumento virtual (\emph{.vi}) que realice lo siguiente:
    \begin{itemize}
        \item Crear una máquina de estados que tenga cuatro estados:
            \begin{itemize}
                \item \emph{stand-by}: espera que se presione un botón y pasa al estado \emph{medición}
                \item \emph{medición}: realiza una medición por segundo por $X$\si{\second} y luego pasa al estado de \emph{espera}
                \item \emph{espera}: espera $Z$\si{\second} y continua al estado de medición a menos que una cantidad de tomas de datos mayor a $Y$ haya sido realizada, en ese caso continua al estado  \emph{fin}
                \item \emph{fin}: termina la ejecución del instrumento virtual.
            \end{itemize}
        \item Guardar los datos tomados en cada estado \emph{medición} en un archivo separado por comas (\emph{.csv}) en \emph{tiempo real}, de forma que la primera columna sea el tiempo y la segunda el valor de la temperatura. El archivo debe cambiar de nombre para cada estado \emph{medición}. 
    \end{itemize}
\item Usando Python realice lo siguiente:
\begin{itemize}
    \item Crear una máquina de estados que tenga cinco estados definidos como:
        \begin{itemize}
            \item \emph{stand-by}: estado que espera que se presione el botón, pasa a estado \emph{hello}
            \item \emph{hello}: estado que concatena letra por letra de la palabra "hello" en un string que inicialmente esta vació, debe aparecer una nueva letra cada segundo, al finalizar la palabra pasa a estado \emph{space}
            \item \emph{space}: estado que concatena un espacio vacío al string y espera por 5 segundos, pasa al estado \emph{world}
            \item \emph{world}: estado que concatena letra por letra la palabra "world" en el string, al finalizar la palabra pasa a estado \emph{end}
            \item \emph{end}: estado que termina la ejecución del programa.
        \end{itemize}
    \item Mostrar como avanza la concatenación del string
    \item Mostrar como avanza el tiempo en segundos en cada estado
\end{itemize}
\item Incluya en su informe lo siguiente:
    \begin{itemize}
        \item Una captura de pantalla de cada estado de la máquina 
        \item Una captura de pantalla del panel frontal
        \item Una familia de gráficas de $T$ que incluya todos los datos de los archivos (\emph{.csv}) realizada a partir de los archivos \emph{.csv}, usando Python. Incluya los títulos de los ejes de forma adecuada. 
        \item El archivo \emph{.py} en el que desarrolló la máquina de estados
    \end{itemize}

\end{enumerate}
